% Generated draft: requires mathematical and source review.
\documentclass[11pt]{article}
\usepackage[margin=25mm]{geometry}
\usepackage{fontspec}
\setmainfont{Latin Modern Roman}
\usepackage{amsmath,amssymb,mathtools,unicode-math}
\setmathfont{Latin Modern Math}
\usepackage{enumitem,needspace,xurl,hyperref}
\setlength{\parindent}{0pt}
\setlength{\parskip}{5pt}
\setlength{\emergencystretch}{3em}
\newcommand{\N}{\mathbb{N}}
\newcommand{\Z}{\mathbb{Z}}
\newcommand{\Q}{\mathbb{Q}}
\newcommand{\R}{\mathbb{R}}
\newcommand{\C}{\mathbb{C}}
\newcommand{\F}{\mathbb{F}}
\newcommand{\A}{\mathbb{A}}
\newcommand{\PP}{\mathbb{P}}
\newcommand{\E}{\mathbb{E}}
\newcommand{\eps}{\varepsilon}
\newcommand{\dd}{\,\mathrm{d}}
\begin{document}
\section*{UnsolvedMath ID: 20002027}
\textbf{Source problem code:} catalogue entry
\subsection{Definitions and mathematical statement}
Problem 8 asks to express the Gauss--Bonnet integrand as a Schouten symmetric-function term plus explicit Weyl corrections, and then explain the relation between Q-curvature and topology. Fix \(n=2m\geq4\), the Laplacian \(\Delta=\operatorname{div}\nabla\), and
\[
 J=\frac{R}{2(n-1)},\qquad
 P=\frac{\operatorname{Ric}-Jg}{n-2},\qquad
 \operatorname{Rm}=W+P\mathbin{\star} g.
\]
Use the Kulkarni--Nomizu convention
\[
 (P\mathbin{\star} g)_{ij}{}^{kl}
 =P_i{}^k\delta_j^l+P_j{}^l\delta_i^k
  -P_i{}^l\delta_j^k-P_j{}^k\delta_i^l.
\]
Let the generalized delta be the determinant of Kronecker deltas and define the standard Euler scalar
\[
 E_{2m}=2^{-m}\delta^{i_1\cdots i_{2m}}_{j_1\cdots j_{2m}}
 \prod_{a=1}^m R_{i_{2a-1}i_{2a}}{}^{j_{2a-1}j_{2a}}.
\]
With these conventions \(E_4=|\operatorname{Rm}|^2-4|\operatorname{Ric}|^2+R^2\), and on a closed manifold
\(\int E_{2m}\,\dd\mathrm{vol}=(4\pi)^m m!\chi(M)\).

The aid's all-even-dimensional finite contraction formula is
\[
 E_{2m}=2^m(m!)^2\sigma_m(P)+\sum_{r=1}^m C_{m,r}(W,P),
\]
where an explicit definition of each correction is
\[
 C_{m,r}=2^{-m}\binom mr
 \delta^{i_1\cdots i_{2m}}_{j_1\cdots j_{2m}}
 \left(\prod_{a=1}^r
 W_{i_{2a-1}i_{2a}}{}^{j_{2a-1}j_{2a}}\right)
 \left(\prod_{a=r+1}^m
 (P\mathbin{\star} g)_{i_{2a-1}i_{2a}}{}^{j_{2a-1}j_{2a}}\right).
\]
Empty products are one. Every correction contains a Weyl factor, and the displayed coefficient of the pure Schouten term fixes its normalization. This is a complete contraction expression, not a claim that every correction has been simplified into a minimal list of independent invariants. For dimension two, where the usual Schouten formula is undefined, the Euler scalar is directly \(E_2=R=2K\).

In dimension four use the full tensor norm \(|W|^2=W_{ijkl}W^{ijkl}\), rather than a curvature-operator norm differing by a fixed factor. The aid gives
\[
 E_4=|W|^2+16\sigma_2(P),\qquad
 Q_4=4\sigma_2(P)-\Delta J,
 \quad J=R/6.
\]
Equivalently \(Q_4=-\frac16\Delta R-\frac12|\operatorname{Ric}|^2+\frac16R^2\). Thus every closed Riemannian four-manifold satisfies
\[
 \int_MQ_4\,\dd\mathrm{vol}
 =8\pi^2\chi(M)-\frac14\int_M|W|^2\,\dd\mathrm{vol}.
\]
In particular total Q is bounded above by \(8\pi^2\chi\), with equality exactly when \(W\) vanishes identically. The closedness assumption removes the Laplacian integral and the Gauss--Bonnet boundary correction.

For the worked family, let \(S^2(a)\) and \(S^2(b)\) denote round spheres of positive Gaussian curvatures \(a,b\), not radii. On their product put \(r=a/b\). Then
\[
 \int Q_4\,\dd\mathrm{vol}
 =\frac{16\pi^2}{3}(4-r-r^{-1}),\qquad
 \int|W|^2\,\dd\mathrm{vol}
 =\frac{64\pi^2}{3}(r+2+r^{-1}).
\]
The Euler characteristic is four. Within this product family total Q is at most \(32\pi^2/3\), with equality precisely at \(a=b\), and tends to \(-\infty\) as \(r\to0\) or \(r\to\infty\). This family maximum is not the unrestricted topological upper bound. The unit round \(S^4\) independently has \(E_4=24\), \(Q_4=6\), \(W=0\), and \(\int Q_4=16\pi^2\), checking the normalization. The aid does not claim a complete all-even-dimensional Q-topology formula or a fully simplified basis for all Weyl corrections.
\subsection{Short English statement}
Expanding the Euler scalar gives an explicit Schouten term and finitely many contractions containing Weyl curvature. In dimension four the resulting integral Q identity separates topology from the Weyl norm. The sphere-product family checks the constants and shows total Q can decrease without bound even at fixed topology.
\subsection{Sources}
\begin{itemize}
\item Permanent UnsolvedMath ID \(20002027\), exact public code catalogue entry; local registry \texttt{aim-geometry-notes.json}, \texttt{notes[364]}, canonical tag problem.
\item AIM workshop ``Conformal structure in geometry, analysis, and physics,'' Problem 8: \url{https://aimath.org/WWN/confstruct/confstruct.pdf}. No PDF was accessed.
\item The supplied verification report says both official PDF and HTML have the typo ``expess,'' the PDF line break explains ``in-volving,'' and \(\sigma_{n/2}(P)\) is the intended notation. The clean-statement field is itself malformed, so the original request and supplied mathematical recovery are distinguished.
\item Research attempt 1, ``Euler integrands, Schouten symmetric functions, and total Q-curvature,'' labelled partial\_result, partially\_solved, proved and low novelty confidence; statement corrected\_verified. It supplies no full primary reference, original literature entry, proposer or year. It explicitly limits the complete Q-topology treatment to dimension four.
\end{itemize}
\end{document}
